\section{Related Work}
\label{sec:related}

\paragraph{Classical numerical methods for MFGs.}
Finite-difference schemes for MFGs were initiated by
\citet{achdou2010mean} and remain the most widely used reference solvers
for low-dimensional problems.
They handle non-separable Hamiltonians but suffer from the curse of
dimensionality, motivating the deep-learning approaches discussed below.
For a recent survey covering both classical and machine-learning methods,
see \citet{hu2024recent}.

\paragraph{GAN-based deep MFG solvers.}
APAC-Net \citep{lin2021apac} and the deep-MFG framework of
\citet{ruthotto2020machine} parameterise the population density by a
generator network and the value function by a critic, exploiting the
primal-dual structure of separable MFGs to obtain a Wasserstein-GAN-like
training objective \citep{arjovsky2017wasserstein,goodfellow2014generative}.
The connection between the saddle-point formulation of separable MFGs
and Wasserstein GANs was made precise by \citet{cao2021connecting}.
Both APAC-Net and \citet{ruthotto2020machine} require separability of $H$,
which is the obstruction addressed by the outer reduction proposed here.

\paragraph{Physics-informed deep MFG solvers.}
\citet{assouline2023deep} (MFDGM) extend the deep-Galerkin method
\citep{sirignano2018dgm,raissi2019physics} to MFGs by minimising the
HJB and Fokker-Planck residuals jointly.
MFDGM handles arbitrary non-separable Hamiltonians but produces only
$L^p$-pointwise approximations and provides no convergence rate.
\Cref{tab:comparison} in \Cref{sec:prelim} summarises the resulting
gap between the GAN and physics-informed families.
\citet{carmona2021deep} provide convergence analysis for related
machine-learning algorithms in the ergodic mean field control setting,
again under separability assumptions.

\paragraph{Fictitious play and policy iteration for MFGs.}
\citet{cardaliaguet2017learning} establish convergence of fictitious play
for separable MFGs under monotonicity, and \citet{perrin2020fictitious}
extend the analysis to a continuous-time learning dynamic.
\citet{lavigne2023policy} propose a policy iteration scheme for
non-separable MFGs that decouples the system into a sequence of linear
PDEs, conceptually parallel to our outer iteration.
Our fixed-point reduction inherits the same decoupling structure, and it can
retain distributional learning whenever the separable inner problem is solved
by a generator-based adversarial method.  The contraction factor \(\kappa_\eps\) in
\Cref{thm:main} gives an explicit exact-outer-map rate in the weak-coupling
regime; in the density-only subcase it reduces to
\(\varepsilon L_R/\lambda\).

\paragraph{Position of our contribution.}
Across the deep-learning MFG literature, no existing method
simultaneously (i) handles non-separable Hamiltonians,
(ii) learns $\rho$ distributionally via a generator,
and (iii) provides Wasserstein convergence with an explicit rate.
The present fixed-point reduction occupies this gap conditionally: the outer
map has an explicit rate in the weak-coupling regime, and distributional
learning is inherited when the separable inner solver is GAN-based.  The clean
threshold \(\varepsilon L_R<\lambda\) is sharp for the density-only linearized
family in \Cref{prop:necessity}.
